\documentclass[11pt]{article}
\usepackage[margin=1in]{geometry}
\usepackage{amsmath,amssymb}
\usepackage{bm}
\usepackage{booktabs}

% --- Matrix / notation macros (consistent with the predecessor description) ---
\newcommand{\rmB}{\mathbf{B}}
\newcommand{\rmX}{\mathbf{X}}
\newcommand{\rmM}{\mathbf{M}}
\newcommand{\rmY}{\mathbf{Y}}
\newcommand{\rmOmega}{\bm{\Omega}}
% If the host document already provides \cref / \autoref (cleveref/hyperref),
% delete the two lines below; they are stubs so this file compiles standalone.
\providecommand{\cref}[1]{Eq.~(\ref{#1})}
\providecommand{\autoref}[1]{\ref{#1}}

\title{The Meta-Simulator}
\author{CogFormer-2}
\date{}

\begin{document}
\maketitle

%============================================================================
% Drop-in replacement for the predecessor's meta-simulator paragraphs.
% Mirrors the current training code path:
%   cogformer/simulator/design_utils.py, coefficient_priors.py,
%   model_family.py, model_class.py;
%   experiments/experiment_1/model_class_cf_train.py + experiments/utils.py
%   (values in Table "tab:meta-sim-hyper").
%============================================================================

To realize the generative hierarchy in \cref{eq:generative}, we build a meta-simulator that treats the parameters of a cognitive model as a generalized linear model (GLM) of the experimental design (see~\autoref{fig:meta-simulator}). A model family $\mathcal{F}$ has $D_{\mathcal{F}}$ intrinsic parameters $\theta_{\mathcal{F}}$; the DDM, for example, has drift rate $v$, threshold $a$, bias $z$, non-decision time $\tau$, variabilities $s_v$ and $s_\tau$, and contamination probability $\pi$. Given a design matrix $\rmX \in \mathbb{R}^{N \times I}$ over $N$ trials, a coefficient matrix $\rmB \in \mathbb{R}^{I \times D_{\mathcal{F}}}$, and a binary mask $\rmM$ of the same shape that marks which coefficients are active, a dataset $\rmY \in \mathbb{R}^{N \times O_{\mathcal{F}}}$ is generated as
\begin{equation}
\label{eq:meta-simulator}
  \Theta = g\!\left(\rmX\,(\rmM \odot \rmB)\right),
  \qquad
  \bm{y}_n \sim (1 - \pi)\, p_{\mathcal{F}}(\bm{y} \mid \bm{\theta}_n)
              + \pi\, \mathcal{U}(\mathcal{Y}_{\mathcal{F}}),
\end{equation}
where row $\bm{\theta}_n$ of $\Theta \in \mathbb{R}^{N \times D_{\mathcal{F}}}$ holds the parameters of trial $n$. The link functions $g$ act column-wise and map each linear predictor to the valid domain of its parameter, using softplus for positive parameters and a scaled sigmoid for bounded ones such as $z$. The second term is a contamination mixture: with probability $\pi$, a trial is replaced by a uniform guess over the response domain, with a response time drawn from $\mathcal{U}(0, 20)$ and a choice drawn uniformly over the alternatives (a uniform angle for the CDM). During training, every dataset draws its own family, design, mask, and coefficients, so that a single network learns to perform inference across all of them.

Intrinsic parameters differ in how they enter the GLM. Core parameters (e.g., $v$, $a$, $z$, $\tau$) are always estimated and may depend on the design. Optional parameters ($s_v$, $s_\tau$) are either estimated or fixed at zero, which removes trial-to-trial variability; a fixed parameter bypasses the link function. The contamination probability $\pi$ is a global parameter: it is estimated once per dataset, never regressed, and drawn directly from $\mathrm{Beta}(1, 50)$ with an identity link. For each dataset, we draw the number of regressors $R \sim \mathcal{U}\{0, \dots, R_{\max}\}$ and two inclusion rates $p_0, p_1 \sim \mathrm{Beta}(2, 2)$. Each optional parameter is included with probability $p_0$, and each included non-global parameter receives a slope on each regressor with probability $p_1$. Two regressors interact on exactly those parameters that both of them carry. Slopes therefore exist only for estimated parameters, and interactions only for main effects that are present. Because $p_0$ and $p_1$ vary across datasets, the network sees designs ranging from intercept-only to fully regressed. During validation, this structure is instead fixed to the benchmark designs in \autoref{app:design-config}.

The first column of $\rmX$ is the intercept, followed by one block of columns per regressor and per interaction. Each regressor is categorical with probability $q_{\mathrm{cat}}$ and continuous otherwise. We draw the categorical regressors first and the continuous regressors conditional on them. This ordering lets us control multicollinearity separately among categorical regressors, between categorical and continuous regressors, and among continuous regressors. A categorical regressor $u_r$ has $k_r \sim \mathcal{U}\{C_{\min}, \dots, C_{\max}\}$ levels with proportions $\bm{p}_r$. These proportions are exactly balanced with probability $q_{\mathrm{bal}}$ and are otherwise drawn from $\mathrm{Dirichlet}(\alpha\,\bm{1})$ with $\alpha \sim \mathrm{LogUniform}(0.5, 50)$, subject to a minimum expected count of two trials per level. Trials are then assigned to combinations of levels according to the joint table
\begin{equation}
\label{eq:joint-category-table}
  \bm{P} = (1 - \lambda) \bigotimes_{r} \bm{p}_r
         + \lambda\, \mathrm{IPF}\!\left(\bm{Q};\, \{\bm{p}_r\}\right),
  \qquad
  \bm{Q} \sim \mathrm{Dirichlet}(\kappa\,\bm{1}),
\end{equation}
which mixes independent factors with a random coupling $\bm{Q}$, rescaled to the same marginals by iterative proportional fitting ($\mathrm{IPF}$). At $\lambda = 0$, the design is an orthogonal factorial. Larger values of $\lambda$ make the categorical regressors increasingly dependent without changing the proportions of any single regressor. Each categorical regressor is dummy-coded against a randomly chosen reference level, so that its levels carry no implied order.

Each continuous regressor combines a component explained by the categorical regressors with a correlated residual:
\begin{equation}
\label{eq:continuous-columns}
  \bm{x}_j = \sqrt{\eta^2}\; \bm{s}_j + \sqrt{1 - \eta^2}\; \bm{\varepsilon}_j,
  \qquad
  \bm{\varepsilon}_n \sim \mathcal{N}(\bm{0}, \rmOmega),
  \quad
  \rmOmega \sim \mathrm{LKJ}(\zeta),
\end{equation}
where $\bm{s}_j$ sums one Gaussian offset per level of each categorical regressor, and both terms are standardized. Thus, $\eta^2$ is the share of a continuous regressor's variance that the categorical regressors explain, and $\rmOmega$ sets the correlations among the continuous regressors. We sample $\rmOmega$ from the LKJ distribution, which places a prior directly on correlation matrices. Its single shape parameter $\zeta$ moves the designs from near-singular ($\zeta < 1$) through uniform ($\zeta = 1$) to near-orthogonal ($\zeta > 1$). Interactions enter $\rmX$ as row-wise products of the two regressor blocks. Because the number of levels and their occupancy vary across datasets, a design column can be empty, for example when an interaction cell receives no trials. The data then carry no information about the corresponding coefficients, so we remove them from $\rmM$.

Intercepts are drawn from the family's prior on the pre-link scale, which is Gaussian for all parameters except $\pi$. Slopes on continuous regressors are drawn from $\mathcal{N}(0, 1)$. For a categorical regressor, we draw one latent mean per level, $\mu_k \sim \mathcal{N}(0, 1)$, and define its contrasts as differences from the reference level; interactions between categorical regressors use the corresponding double difference:
\begin{equation}
\label{eq:cell-means}
  \beta_k = \mu_k - \mu_1,
  \qquad
  \beta_{kl} = \mu_{kl} - \mu_{k1} - \mu_{1l} + \mu_{11}.
\end{equation}
Independent contrasts would place the reference level a priori closer to every other level. This cell-means prior instead treats all levels symmetrically, so the choice of reference level does not affect the prior.

Finally, model families share some parameters but not all. We therefore embed every family in a common parameter space of dimension $D_{\mathcal{C}}$, the ordered union of all intrinsic parameters in the model class $\mathcal{C}$, and parameters absent from the sampled family remain zero and masked. In the same way, $\rmX$, $\rmM$, and $\rmB$ are zero-padded to $I_{\max}$ columns, the size of a design with $R_{\max}$ regressors of $C_{\max}$ levels and all of their pairwise interactions. The number of trials $N$ is shared within a batch, whereas all other quantities are drawn for each dataset. \autoref{app:meta-sim-settings} lists the settings used for training.

%============================================================================
% APPENDIX BLOCK: move into the paper's appendix (e.g., next to
% app:design-config). Drop the \appendix line if the host document already
% has one before this point.
%============================================================================
\appendix
\section{Meta-simulator settings}
\label{app:meta-sim-settings}

\autoref{tab:meta-sim-hyper} lists the distributions and constants of the meta-simulator used to train the model-class network. Symbols follow the main text.

\begin{table}[h]
\centering
\caption{Meta-simulator settings used for training.}
\label{tab:meta-sim-hyper}
\begin{tabular}{lll}
  \toprule
  Quantity & Symbol & Value \\
  \midrule
  Model families               & $\mathcal{C}$          & DDM, RDM, CDM (uniform) \\
  Shared parameter space       & $D_{\mathcal{C}}$      & $9$ ($v, a, z, \tau, s_v, s_\tau, \pi, v_{\mathrm{diff}}, v_\theta$) \\
  Trials per dataset           & $N$                    & $\mathcal{U}\{200, \dots, 1000\}$ \\
  Regressors                   & $R$                    & $\mathcal{U}\{0, \dots, 4\}$ \\
  Levels per categorical regressor & $k_r$              & $\mathcal{U}\{2, 3, 4\}$ \\
  Design capacity              & $I_{\max}$             & $67$ \\
  Inclusion rates              & $p_0, p_1$             & $\mathrm{Beta}(2, 2)$ \\
  Contamination                & $\pi$                  & $\mathrm{Beta}(1, 50)$ \\
  \midrule
  Categorical probability      & $q_{\mathrm{cat}}$     & $0.5$ \\
  Balanced-level probability   & $q_{\mathrm{bal}}$     & $0.3$ \\
  Level concentration          & $\alpha$               & $\mathrm{LogUniform}(0.5, 50)$ \\
  Categorical association      & $\lambda$              & $0$ w.p.\ $0.5$, else $\mathcal{U}(0, 1)$ \\
  Coupling concentration       & $\kappa$               & $\mathrm{LogUniform}(0.3, 10)$ \\
  Categorical--continuous association & $\eta^2$        & $\min\{\mathrm{Beta}(1.2, 3),\, 0.85\}$ \\
  LKJ shape                    & $\zeta$                & $\mathcal{U}(0.5, 5)$ \\
  \bottomrule
\end{tabular}
\end{table}

\end{document}
